\section{Discussion}
\label{sec:discussion}

\TODO{This section is deliberately left as a placeholder. The final discussion and scope statement depend on the outcome of the confound-attribution audit's generalization to GSE146889, CPTAC-CCRCC, and TCGA-LUAD (\Cref{sec:methods-confound}, \Cref{sec:results-confound}), which is in progress and has not yet reported back at the time of writing. Do not draft a synthesis here that assumes a particular outcome for those three audits --- the pilot cohort's finding (confound-independence of the $H_1$ signal, scoped strictly to GSE81089) must not be extrapolated to a general or cross-cohort claim until it is either replicated or refuted in each of the other three cohorts. Note in particular that GSE146889's own cross-check (\Cref{sec:results-confound-crossdataset}) already shows a structural difference from the pilot --- a significantly tumor-enriched dominant $H_1$ loop --- that is a specific, concrete reason the pilot's confound-independence result might not transfer there; the synthesis phase should treat this as a live possibility to be resolved by the pending audit, not dismissed in advance.}

\TODO{Once the three remaining confound audits complete, this section should address at minimum: (i) what the complete cross-cohort pattern implies about whether the pilot's confound-independence finding is a general property of the pipeline or a property specific to the pilot cohort's biology/data-collection process; (ii) how the ablation sweep's identified failure boundaries (PC=10; HVG=4000/PC=100) interact with the confound-audit's scope --- e.g., whether a dataset that fails the ablation boundary conditions should be interpreted differently in the confound audit; (iii) an honest statement of what this program does and does not establish about the general relationship between PCA and persistent-homology-based topological claims in omics data, calibrated against the actual evidence gathered rather than the most publishable-sounding interpretation of it; (iv) limitations, including the reliance on a single primary test statistic (max-$H_1$ persistence) rather than a full persistence-diagram distance, the restriction to five cohorts across a still-limited set of tissue types, and the fact that the ablation sweep was run only on the pilot cohort and has not itself been cross-validated on the replication cohorts.}
